\chapter{Interstitial-Condensate Constraint and Defect Impedance}
\label{app:condensate_impedance}

\begin{abstract}
\noindent This appendix establishes the mechanical response of the interstitial condensate within the active FCC vacuum. It defines the incrementally incompressible coherent branch, separating the singular bulk constraint from the finite gravitational yield pressure. It then formulates the localized relative-flow action for a topological defect, constructing the mode-resolved added-inertia functional and defining the microscopic impedance projection as a boundary-value functional as a mode-weighted inverse penetration length. Canonical defect-flow solutions and the numerical evaluation of the projection remain open boundary-value problems.
\end{abstract}

\section{Scope and Constitutive Setting}

The active vacuum is a saturated dual-phase medium composed of the solid FCC scaffold and the connected interstitial condensate. The primitive-cell volume is partitioned according to
\begin{equation}
    \beta_{\mathrm{geo}} = \frac{\pi}{3\sqrt{2}},
\end{equation}
and
\begin{equation}
    \beta_{\mathrm{IC}} = 1-\beta_{\mathrm{geo}},
\end{equation}
so that
\begin{equation}
    \beta_{\mathrm{geo}}+\beta_{\mathrm{IC}}=1.
\end{equation}

The FCC nodes occupy the close-packed solid fraction, while the interstitial condensate continuously fills the complementary octahedral and tetrahedral regions. No unoccupied third phase is introduced within the active FCC cell.

The analysis is restricted to the active phase,
\begin{equation}
    0 \leq \alpha_{\mathrm{G}} < 1,
\end{equation}
and to perturbations about a gravitationally preloaded background that is quasistatic over the duration of the local disturbance,
\begin{equation}
    \dot{\alpha}_{\mathrm{G}} \approx 0.
\end{equation}

The loading parameter $\alpha_{\mathrm{G}}$ specifies the background state about which the incremental dynamics are constructed. It is not itself treated as a propagating perturbation in the coherent branch developed below.

The interstitial condensate performs three mechanically distinct functions. It participates in the inertia of coherent scaffold--condensate motion, supports the incremental pressure constraint of the saturated active cell, and develops relative-flow impedance around localized topological defects. These functions must not be combined into a single undifferentiated density or kinetic term.

The coherent periodic branch is characterized by entrainment of the condensate with the FCC scaffold. In this branch, the condensate may be displaced, but it does not move relative to the scaffold. The localized defect branch is characterized by nonzero relative condensate motion around a moving or internally reconfiguring defect. The two branches correspond to different reductions of the same underlying dual-phase medium.

Let $\mathbf{u}(\mathbf{x},t)$ denote the homogenized FCC scaffold displacement and $\mathbf{w}(\mathbf{x},t)$ the interstitial-condensate displacement. The relative condensate velocity is
\begin{equation}
    \mathbf{v}_{\mathrm{rel}} = \dot{\mathbf{w}} - \dot{\mathbf{u}}.
\end{equation}

For the coherent active-vacuum branch, $\mathbf{v}_{\mathrm{rel}} = \mathbf{0}$, whereas for a localized defect response, $\mathbf{v}_{\mathrm{rel}} \neq \mathbf{0}$ within the condensate domain influenced by the defect.

The present appendix first formulates the coherent condensate constraint and then develops the localized relative-flow action. The coherent sector determines the admissible translational acoustic polarizations and the incremental pressure field. The localized sector determines the mode-dependent hydrodynamic impedance entering the defect mass law.

The finite gravitational loading of the active vacuum must also be distinguished from the incremental constitutive response. The background may sustain a finite pressure and may approach the FCC saturation threshold as $\alpha_{\mathrm{G}} \rightarrow 1^{-}$. The coherent perturbation about that background may nevertheless be incrementally volume preserving. Thus,
\begin{equation}
    \text{finite background loading} \neq \text{incremental volumetric wave}.
\end{equation}

The pressure constraint derived below applies to the low-energy coherent branch within the active phase. It does not prohibit quasistatic background compaction, nonlinear lattice rearrangement, or the FCC-to-SC phase transition at saturation.

\section{Coherent Scaffold--Condensate Kinematics}

In the entrained coherent branch, $\mathbf{w}=\mathbf{u}$ up to an irrelevant time-independent reference displacement. Consequently, $\dot{\mathbf{w}}=\dot{\mathbf{u}}$, and $\mathbf{v}_{\mathrm{rel}}=\mathbf{0}$. The FCC scaffold and the interstitial condensate therefore share one coherent displacement field, $\mathbf{u}(\mathbf{x},t)$.

The associated infinitesimal strain tensor is
\begin{equation}
    \varepsilon_{ij} = \frac{1}{2} \left( \partial_i u_j + \partial_j u_i \right).
\end{equation}

Its volumetric component is $\varepsilon_V = \operatorname{tr}\boldsymbol{\varepsilon} = \nabla\cdot\mathbf{u}$, and its deviatoric part is $\varepsilon_{ij}^{\mathrm{dev}} = \varepsilon_{ij} - \frac{1}{3}\varepsilon_V\delta_{ij}$.

Let $\rho_{\mathrm{eff}}^{(0)}$ denote the effective unstrained inertia density of the selected coherent continuum description. The same formalism applies whether this coefficient is interpreted at the primitive-cell microscopic scale or at a coarse-grained macroscopic scale, provided it is paired consistently with the corresponding shear coefficient.

The gravitational endpoint-transfer constraint modifies the completed-coordinate inertia according to
\begin{equation}
    \rho_{\mathrm{eff}}(\alpha_{\mathrm{G}}) = \frac{\rho_{\mathrm{eff}}^{(0)}}{(1-\alpha_{\mathrm{G}})^4}.
\end{equation}

The coherent kinetic-energy density is therefore
\begin{equation}
    \mathcal{T}_{\mathrm{coh}} = \frac{1}{2} \rho_{\mathrm{eff}}(\alpha_{\mathrm{G}}) \dot{\mathbf{u}}\cdot\dot{\mathbf{u}}.
\end{equation}

The coefficient $\rho_{\mathrm{eff}}(\alpha_{\mathrm{G}})$ already represents the inertia of the entrained coherent transfer coordinate. The localized condensate impedance density introduced later must not be added independently to this coherent kinetic term. Such an addition would count the condensate response twice: once as entrained coherent inertia and again as localized relative-flow inertia.

Let $K_{\mathrm{sh}}^{(0)}$ denote the shear coefficient paired with $\rho_{\mathrm{eff}}^{(0)}$ at the same effective scale. The unstrained transverse normalization is
\begin{equation}
    K_{\mathrm{sh}}^{(0)} = \rho_{\mathrm{eff}}^{(0)}c_{\mathrm{E}}^2.
\end{equation}

At the retained active-phase tangent order, gravitational loading modifies the completed-coordinate inertia but does not continuously soften the shear coefficient:
\begin{equation}
    K_{\mathrm{sh}}(\alpha_{\mathrm{G}}) = K_{\mathrm{sh}}^{(0)}, \qquad 0 \leq \alpha_{\mathrm{G}} < 1.
\end{equation}

Before imposing the coherent volume constraint, introduce a finite volumetric regularization coefficient $B_{\mathrm{IC}}$. The regularized elastic-energy density is
\begin{equation}
    \mathcal{W}_B = K_{\mathrm{sh}}^{(0)} \varepsilon_{ij}\varepsilon_{ij} + \frac{1}{2} B_{\mathrm{IC}} \varepsilon_V^2.
\end{equation}

The corresponding isotropic tangent tensor is
\begin{equation}
    C_{ijkl}^{(B)} = B_{\mathrm{IC}}\delta_{ij}\delta_{kl} + K_{\mathrm{sh}}^{(0)} \left( \delta_{ik}\delta_{jl} + \delta_{il}\delta_{jk} \right).
\end{equation}

The corresponding long-wavelength acoustic tensor is obtained by projection along the propagation direction. Define
\begin{equation}
    \hat{\mathbf s}
    =
    \frac{\mathbf k}{|\mathbf k|}.
\end{equation}
The acoustic tensor is
\begin{equation}
    \Gamma_{ij}
    =
    C_{ikjl}^{(B)}s_k s_l.
\end{equation}
Using the regularized isotropic tangent tensor above, this becomes
\begin{equation}
    \Gamma_{ij}
    =
    K_{\mathrm{sh}}^{(0)}\delta_{ij}
    +
    \left(
    B_{\mathrm{IC}}
    +
    K_{\mathrm{sh}}^{(0)}
    \right)
    s_i s_j.
\end{equation}

For a transverse polarization satisfying
\begin{equation}
    \hat{\mathbf s}\cdot\mathbf e_T=0,
\end{equation}
one obtains
\begin{equation}
    \Gamma_{ij}e_{T,j}
    =
    K_{\mathrm{sh}}^{(0)}e_{T,i}.
\end{equation}
For a longitudinal polarization satisfying
\begin{equation}
    \mathbf e_L\parallel\hat{\mathbf s},
\end{equation}
one obtains
\begin{equation}
    \Gamma_{ij}e_{L,j}
    =
    \left(
    B_{\mathrm{IC}}
    +
    2K_{\mathrm{sh}}^{(0)}
    \right)
    e_{L,i}.
\end{equation}

Thus the finite-compressibility regularization contains two transverse acoustic eigenvalues
\begin{equation}
    \lambda_T
    =
    K_{\mathrm{sh}}^{(0)},
\end{equation}
and one longitudinal acoustic eigenvalue
\begin{equation}
    \lambda_L
    =
    B_{\mathrm{IC}}
    +
    2K_{\mathrm{sh}}^{(0)}.
\end{equation}
This acoustic-tensor projection gives the direct tensor bridge between the regularized constitutive law and the transverse and longitudinal dispersion relations derived below.

The total bulk coefficient associated with this finite regularization is
\begin{equation}
    K_{\mathrm{bulk}}^{(B)} = B_{\mathrm{IC}} + \frac{2}{3} K_{\mathrm{sh}}^{(0)}.
\end{equation}

The coefficient $B_{\mathrm{IC}}$ is introduced only to expose the structure of a hypothetical ordinary longitudinal acoustic branch. It is not yet identified with the finite FCC saturation-pressure scale. The tangent volume constraint and the finite phase-transition pressure are distinct constitutive quantities.

\section{Finite-Compressibility Regularization}

The quadratic action of the regularized coherent branch is
\begin{equation}
    S_B^{(2)} = \int \left[ \frac{1}{2} \rho_{\mathrm{eff}}(\alpha_{\mathrm{G}}) \dot{u}_i\dot{u}_i - K_{\mathrm{sh}}^{(0)} \varepsilon_{ij}\varepsilon_{ij} - \frac{1}{2} B_{\mathrm{IC}} \varepsilon_V^2 \right] \mathrm{d}^3x \, \mathrm{d}t.
\end{equation}

Variation with respect to the displacement field gives the incremental stress
\begin{equation}
    \sigma_{ij} = 2K_{\mathrm{sh}}^{(0)} \varepsilon_{ij} + B_{\mathrm{IC}} \varepsilon_V\delta_{ij}.
\end{equation}

The equation of motion is $\rho_{\mathrm{eff}}(\alpha_{\mathrm{G}}) \ddot{u}_i = \partial_j\sigma_{ij}$. Using the linear strain definition, $\partial_j\varepsilon_{ij} = \frac{1}{2} \left[ \nabla^2u_i + \partial_i(\nabla\cdot\mathbf{u}) \right]$, the vector equation becomes
\begin{equation}
    \rho_{\mathrm{eff}}(\alpha_{\mathrm{G}}) \ddot{\mathbf{u}} = K_{\mathrm{sh}}^{(0)} \nabla^2\mathbf{u} + \left( B_{\mathrm{IC}} + K_{\mathrm{sh}}^{(0)} \right) \nabla(\nabla\cdot\mathbf{u}).
\end{equation}

Consider a plane-wave perturbation,
\begin{equation}
    \mathbf{u} = \mathbf{e} \exp\left[ i(\mathbf{k}\cdot\mathbf{x}-\omega t) \right].
\end{equation}

For a transverse polarization ($\mathbf{k}\cdot\mathbf{e}_T=0$), the dispersion relation is
\begin{equation}
    \omega_T^2 = \frac{ K_{\mathrm{sh}}^{(0)} }{ \rho_{\mathrm{eff}}(\alpha_{\mathrm{G}}) } k^2 \quad \implies \quad c_T^2(\alpha_{\mathrm{G}}) = \frac{ K_{\mathrm{sh}}^{(0)} }{ \rho_{\mathrm{eff}}(\alpha_{\mathrm{G}}) }.
\end{equation}

For a longitudinal polarization ($\mathbf{e}_L\parallel\mathbf{k}$), the dispersion relation is
\begin{equation}
    \omega_L^2 = \frac{ B_{\mathrm{IC}} + 2K_{\mathrm{sh}}^{(0)} }{ \rho_{\mathrm{eff}}(\alpha_{\mathrm{G}}) } k^2 \quad \implies \quad c_L^2(\alpha_{\mathrm{G}}) = \frac{ B_{\mathrm{IC}} + 2K_{\mathrm{sh}}^{(0)} }{ \rho_{\mathrm{eff}}(\alpha_{\mathrm{G}}) }.
\end{equation}

In the unstrained active vacuum ($\alpha_{\mathrm{G}}=0$), $\rho_{\mathrm{eff}}(0) = \rho_{\mathrm{eff}}^{(0)}$. Using $K_{\mathrm{sh}}^{(0)} = \rho_{\mathrm{eff}}^{(0)}c_{\mathrm{E}}^2$, the transverse speed becomes $c_T(0)=c_{\mathrm{E}}$. The hypothetical longitudinal speed is
\begin{equation}
    c_L^2(0) = c_{\mathrm{E}}^2 \left( \frac{B_{\mathrm{IC}}}{K_{\mathrm{sh}}^{(0)}} + 2 \right).
\end{equation}

Mechanical stability of the finite-compressibility regularization requires $K_{\mathrm{bulk}}^{(B)}>0$, or
\begin{equation}
    B_{\mathrm{IC}} + \frac{2}{3} K_{\mathrm{sh}}^{(0)} > 0 \quad \implies \quad \frac{B_{\mathrm{IC}}}{K_{\mathrm{sh}}^{(0)}} > -\frac{2}{3}.
\end{equation}

It follows that every mechanically stable finite ordinary longitudinal branch satisfies
\begin{equation}
    c_L^2(0) > \frac{4}{3} c_{\mathrm{E}}^2,
\end{equation}
and hence
\begin{equation}
    \boxed{ c_L(0) > \sqrt{\frac{4}{3}} \, c_{\mathrm{E}} }.
\end{equation}

For a nonnegative regularizing compression coefficient ($B_{\mathrm{IC}}\geq0$), the stronger bound is $c_L(0)\geq\sqrt{2}c_{\mathrm{E}}$.

Accordingly, the following three conditions cannot be imposed simultaneously within an ordinary finite isotropic elastic model:
\begin{itemize}
    \item $c_T=c_{\mathrm{E}}$,
    \item $c_{\mathrm{E}}$ is the active-phase propagation ceiling,
    \item $B_{\mathrm{IC}}$ produces an independent stable finite-frequency longitudinal acoustic branch.
\end{itemize}

The finite-$B_{\mathrm{IC}}$ theory therefore functions as a diagnostic regularization. It demonstrates that an ordinary scalar acoustic polarization is incompatible with the selected kinematic structure of the active Soft Floor once the transverse branch is normalized to $c_{\mathrm{E}}$.

This incompatibility does not uniquely derive the incompressible limit from stability alone. Rather, it motivates the constitutive selection of a reduced coherent branch in which local volume fluctuations are excluded from the finite-energy acoustic spectrum.

\section{Incrementally Incompressible Coherent Branch}

The coherent active-FCC branch is modeled by the singular reduction
\begin{equation}
    \boxed{ B_{\mathrm{IC}}\rightarrow\infty }.
\end{equation}

For a finite-energy perturbation, $\frac{1}{2} B_{\mathrm{IC}} \varepsilon_V^2 < \infty$. Consequently, $\varepsilon_V\rightarrow0$, and the admissible coherent displacement satisfies
\begin{equation}
    \boxed{ \nabla\cdot\mathbf{u}=0 }.
\end{equation}

This is an incremental constraint about a prescribed gravitationally loaded background. It does not require the background FCC cell to possess zero finite compaction, nor does it prohibit the approach to the nonlinear saturation boundary. The background state is encoded by $\alpha_{\mathrm{G}}$; the coherent perturbation about that state is volume preserving.

Introduce the incremental condensate pressure $\delta p_{\mathrm{IC}}(\mathbf{x},t)$ as a Lagrange multiplier enforcing the volume constraint. The reduced coherent action is
\begin{equation}
    S_{\mathrm{coh}}^{(2)} = \int \left[ \frac{1}{2} \rho_{\mathrm{eff}}(\alpha_{\mathrm{G}}) \dot{u}_i\dot{u}_i - K_{\mathrm{sh}}^{(0)} \varepsilon_{ij}\varepsilon_{ij} + \delta p_{\mathrm{IC}} \nabla\cdot\mathbf{u} \right] \mathrm{d}^3x \, \mathrm{d}t.
\end{equation}

The same coherent volume constraint is also obtained from condensate constituent conservation. Let
\begin{equation}
    \varrho_{\mathrm{IC}}
\end{equation}
denote the intrinsic constituent density of the interstitial condensate. This symbol is distinct from the projected defect impedance density introduced later. The local conservation law is
\begin{equation}
    \frac{D\varrho_{\mathrm{IC}}}{Dt}
    +
    \varrho_{\mathrm{IC}}
    \nabla\cdot\mathbf v_{\mathrm{IC}}
    =
    0.
\end{equation}
In the coherent entrained branch,
\begin{equation}
    \mathbf v_{\mathrm{IC}}
    =
    \dot{\mathbf u}.
\end{equation}
If the intrinsic constituent density is fixed on the selected incompressible branch,
\begin{equation}
    \frac{D\varrho_{\mathrm{IC}}}{Dt}
    =
    0,
\end{equation}
then the conservation law gives
\begin{equation}
    \nabla\cdot\dot{\mathbf u}
    =
    0.
\end{equation}
For an initially volume-preserving reference state, this integrates to
\begin{equation}
    \nabla\cdot\mathbf u
    =
    0.
\end{equation}

Thus the coherent constraint follows consistently from both the singular tangent limit
\begin{equation}
    B_{\mathrm{IC}}\rightarrow\infty
\end{equation}
and intrinsic condensate constituent conservation. The notation \(\varrho_{\mathrm{IC}}\) is used here to avoid identifying the intrinsic constituent density with the mode-resolved defect impedance density \(\rho_{\mathrm{IC}}^{(\mathrm{imp})}\).

Variation with respect to $\delta p_{\mathrm{IC}}$ gives $\nabla\cdot\mathbf{u}=0$. Variation with respect to $\mathbf{u}$ gives
\begin{equation}
    \rho_{\mathrm{eff}}(\alpha_{\mathrm{G}}) \ddot{u}_i = \partial_j\delta\sigma_{ij},
\end{equation}
where $\delta\sigma_{ij} = -\delta p_{\mathrm{IC}}\delta_{ij} + 2K_{\mathrm{sh}}^{(0)} \varepsilon_{ij}$.

Using the incompressibility condition, the equation of motion becomes
\begin{equation}
    \rho_{\mathrm{eff}}(\alpha_{\mathrm{G}}) \ddot{\mathbf{u}} = K_{\mathrm{sh}}^{(0)} \nabla^2\mathbf{u} - \nabla\delta p_{\mathrm{IC}}.
\end{equation}

In a source-free homogeneous region, taking the divergence gives $\nabla^2\delta p_{\mathrm{IC}}=0$. In the presence of an applied incremental force density $\mathbf{f}$,
\begin{equation}
    \rho_{\mathrm{eff}}(\alpha_{\mathrm{G}}) \ddot{\mathbf{u}} = K_{\mathrm{sh}}^{(0)} \nabla^2\mathbf{u} - \nabla\delta p_{\mathrm{IC}} + \mathbf{f},
\end{equation}
and the pressure satisfies
\begin{equation}
    \boxed{ \nabla^2\delta p_{\mathrm{IC}} = \nabla\cdot\mathbf{f} }.
\end{equation}

The elliptic pressure equation is the compatibility condition of the reduced constrained branch. It does not define an independently propagating field and does not represent an additional channel of instantaneous physical information. The pressure is determined jointly by the applied loading, the incompressibility constraint, and the boundary conditions.

For a plane-wave perturbation, the constraint requires $\mathbf{k}\cdot\mathbf{e}=0$. A longitudinal translational polarization is therefore excluded from the admissible finite-energy configuration space of the reduced coherent branch.

The formal longitudinal frequency of the finite-$B_{\mathrm{IC}}$ regularization satisfies
\begin{equation}
    \omega_L^2 = \frac{ B_{\mathrm{IC}} + 2K_{\mathrm{sh}}^{(0)} }{ \rho_{\mathrm{eff}}(\alpha_{\mathrm{G}}) } k^2.
\end{equation}

For fixed nonzero $k$, $\omega_L\rightarrow\infty$ as $B_{\mathrm{IC}}\rightarrow\infty$. The longitudinal translational excitation is therefore removed from the finite-frequency coherent spectrum. It is not converted into a superluminal propagating signal.

For either admissible transverse polarization, the pressure term is orthogonal to the polarization vector, and the dispersion relation is
\begin{equation}
    \omega_T^2 = \frac{ K_{\mathrm{sh}}^{(0)} }{ \rho_{\mathrm{eff}}(\alpha_{\mathrm{G}}) } k^2 + \mathcal{O}(k^4d_{\mathrm{nn}}^2).
\end{equation}

Using $K_{\mathrm{sh}}^{(0)} = \rho_{\mathrm{eff}}^{(0)}c_{\mathrm{E}}^2$ and $\rho_{\mathrm{eff}}(\alpha_{\mathrm{G}}) = \rho_{\mathrm{eff}}^{(0)} / (1-\alpha_{\mathrm{G}})^4$, one obtains
\begin{equation}
    \omega_T^2 = c_{\mathrm{E}}^2 (1-\alpha_{\mathrm{G}})^4 k^2 + \mathcal{O}(k^4d_{\mathrm{nn}}^2).
\end{equation}

Hence, the propagation speed is
\begin{equation}
    \boxed{ c_P = c_{\mathrm{E}} (1-\alpha_{\mathrm{G}})^2 }.
\end{equation}

The translational acoustic sector of the coherent active vacuum therefore contains two finite-frequency transverse polarizations. This statement does not classify the independent dynamical status of internal phase, microrotational, torsional, or localized defect modes.

\section{Static Longitudinal Pressure and the Finite FCC Yield Scale}

Incremental incompressibility does not imply the absence of hydrostatic stress. The interstitial-condensate pressure is decomposed into a quasistatic background contribution and an incremental reaction field,
\begin{equation}
    p_{\mathrm{IC}}(\mathbf{x},t) = \overline{p}_{\mathrm{IC}}(\mathbf{x}) + \delta p_{\mathrm{IC}}(\mathbf{x},t).
\end{equation}

The background pressure $\overline{p}_{\mathrm{IC}}$ supports the prescribed gravitational preload. The incremental pressure $\delta p_{\mathrm{IC}}$ enforces the volume constraint of the coherent perturbation.

The local gravitational loading is represented by
\begin{equation}
    \tau_L = \alpha_{\mathrm{G}}\tau_{\mathrm{D}},
\end{equation}
where $\tau_{\mathrm{D}}$ is the active Soft Floor surface-tension scale. The corresponding volumetric pressure scale is obtained through the node-radius projection,
\begin{equation}
    \tau_{\mathrm{D,bulk}} \equiv \frac{\tau_{\mathrm{D}}}{r_{\mathrm{D}}}.
\end{equation}

The quantity $\tau_{\mathrm{D,bulk}}$ has dimensions of pressure and is interpreted as the finite gravitational yield pressure of the active FCC phase. It is not the tangent bulk modulus of an ordinary compressible acoustic medium.

The quasistatic condensate pressure is therefore
\begin{equation}
    \overline{p}_{\mathrm{IC}} = \frac{\tau_L}{r_{\mathrm{D}}},
\end{equation}
and hence
\begin{equation}
    \overline{p}_{\mathrm{IC}} = \alpha_{\mathrm{G}}\tau_{\mathrm{D,bulk}}.
\end{equation}

Equivalently,
\begin{equation}
    \alpha_{\mathrm{G}} = \frac{\overline{p}_{\mathrm{IC}}}{\tau_{\mathrm{D,bulk}}}.
\end{equation}

Within the active phase,
\begin{equation}
    0 \leq \overline{p}_{\mathrm{IC}} < \tau_{\mathrm{D,bulk}}.
\end{equation}

At zero gravitational loading, $\alpha_{\mathrm{G}}=0$, and therefore $\overline{p}_{\mathrm{IC}}=0$.

In the active-side saturation limit, $\alpha_{\mathrm{G}}\rightarrow1^{-}$, the background pressure approaches
\begin{equation}
    \overline{p}_{\mathrm{IC}} \rightarrow \tau_{\mathrm{D,bulk}}.
\end{equation}

At equality,
\begin{equation}
    \overline{p}_{\mathrm{IC}} = \tau_{\mathrm{D,bulk}},
\end{equation}
the active FCC constitutive coordinates terminate. The local structure enters the saturated phase-transition regime and can no longer be represented as a linearly perturbed active FCC cell.

The tangent incompressibility condition and the finite saturation pressure perform distinct constitutive roles. The singular limit
\begin{equation}
    B_{\mathrm{IC}}\rightarrow\infty
\end{equation}
excludes finite coherent volume strain,
\begin{equation}
    \nabla\cdot\mathbf{u}=0,
\end{equation}
whereas
\begin{equation}
    \tau_{\mathrm{D,bulk}}<\infty
\end{equation}
sets the maximum quasistatic pressure sustainable by the active FCC phase. Consequently,
\begin{equation}
    \boxed{ B_{\mathrm{IC}} \neq \tau_{\mathrm{D,bulk}} }.
\end{equation}
The former is a tangent constraint parameter; the latter is a finite yield-pressure scale.

The total incremental stress about the gravitationally preloaded state is
\begin{equation}
    \delta\sigma_{ij} = -\delta p_{\mathrm{IC}}\delta_{ij} + 2K_{\mathrm{sh}}^{(0)} \varepsilon_{ij}.
\end{equation}

The corresponding total stress may be represented as
\begin{equation}
    \sigma_{ij} = -\overline{p}_{\mathrm{IC}}\delta_{ij} - \delta p_{\mathrm{IC}}\delta_{ij} + 2K_{\mathrm{sh}}^{(0)} \varepsilon_{ij},
\end{equation}
with the background term satisfying the quasistatic equilibrium equations independently of the incremental disturbance.

The background pressure determines the available endpoint-transfer capacity. Since
\begin{equation}
    \alpha_{\mathrm{G}} = \frac{\overline{p}_{\mathrm{IC}}}{\tau_{\mathrm{D,bulk}}},
\end{equation}
the locally uniform endpoint-transfer factor becomes
\begin{equation}
    \mathcal{C} = \left( 1- \frac{\overline{p}_{\mathrm{IC}}}{\tau_{\mathrm{D,bulk}}} \right)^2.
\end{equation}

The gravitationally referred coherent inertia is therefore
\begin{equation}
    \rho_{\mathrm{eff}} = \frac{\rho_{\mathrm{eff}}^{(0)}}{ \left( 1- \overline{p}_{\mathrm{IC}} / \tau_{\mathrm{D,bulk}} \right)^4 }.
\end{equation}

The transverse propagation speed follows as
\begin{equation}
    \boxed{ c_P = c_{\mathrm{E}} \left( 1- \frac{\overline{p}_{\mathrm{IC}}}{\tau_{\mathrm{D,bulk}}} \right)^2 }.
\end{equation}

The static longitudinal pressure and the dynamic transverse propagation law are thereby coupled through the occupied active capacity, while remaining mechanically distinct sectors.

\section{Phase-Sector Status and Translational Acoustic Mode Count}

The objective contact action developed in Appendix \ref{app:fcc_microphysics} contains the internal phase difference
\begin{equation}
    \Delta\varphi_{ij} = \varphi_j-\varphi_i,
\end{equation}
with local phase energy
\begin{equation}
    \mathcal{U}_{\varphi}^{(ij)} = \kappa_\varphi r_{\mathrm{D}}^2 \left[ 1-\cos(\Delta\varphi_{ij}) \right].
\end{equation}

In the harmonic regime,
\begin{equation}
    \mathcal{U}_{\varphi}^{(ij)} = \frac{1}{2} \kappa_\varphi r_{\mathrm{D}}^2 (\Delta\varphi_{ij})^2 + \mathcal{O}(\Delta\varphi_{ij}^4).
\end{equation}

Because the contact energy depends only on phase differences, a uniform phase displacement,
\begin{equation}
    \varphi_i \rightarrow \varphi_i+\varphi_0,
\end{equation}
leaves every bond energy unchanged. The corresponding uniform phase direction is therefore a null direction of the phase-difference potential.

This property does not by itself determine whether the internal phase constitutes an independent propagating bulk degree of freedom. Such a determination requires specification of the phase kinetic sector.

If the phase possesses an independent kinetic coefficient $\mathcal{I}_\varphi$, a continuum phase action may take the local form
\begin{equation}
    S_\varphi = \int \left[ \frac{1}{2} \mathcal{I}_\varphi \dot{\varphi}^2 - \frac{1}{2} K_\varphi |\nabla\varphi|^2 - U_{\mathrm{lock}}(\varphi) \right] \mathrm{d}^3x \, \mathrm{d}t.
\end{equation}

When $U_{\mathrm{lock}}(\varphi)=0$, the linearized phase branch satisfies
\begin{equation}
    \omega_\varphi^2 = \frac{K_\varphi}{\mathcal{I}_\varphi} k^2.
\end{equation}

When a local locking curvature is present,
\begin{equation}
    m_\varphi^2 \equiv \left. \frac{\mathrm{d}^2U_{\mathrm{lock}}}{\mathrm{d}\varphi^2} \right|_{\varphi=\varphi_0},
\end{equation}
the branch becomes
\begin{equation}
    \omega_\varphi^2 = \frac{K_\varphi}{\mathcal{I}_\varphi} k^2 + \frac{m_\varphi^2}{\mathcal{I}_\varphi}.
\end{equation}

If no independent phase inertia is assigned, $\varphi$ functions instead as a constrained, gauge-redundant, or defect-confined coordinate and does not define an additional bulk propagating mode.

The present appendix does not select among these possibilities. Its coherent-mode statement is restricted to the translational acoustic sector. Within that sector, the incompressibility constraint
\begin{equation}
    \nabla\cdot\mathbf{u}=0
\end{equation}
removes the longitudinal translational polarization and leaves two transverse translational acoustic polarizations.

Accordingly, the statement that \textit{the coherent translational acoustic sector contains two finite-frequency transverse polarizations} does not imply that the complete microscopic spectrum contains only two dynamical branches. Microrotational, torsional, internal-phase, chiral, and localized-defect branches remain outside the translational acoustic mode count.

For either translational transverse polarization,
\begin{equation}
    \mathbf{k}\cdot\mathbf{e}_T=0,
\end{equation}
and
\begin{equation}
    \omega_T^2 = c_{\mathrm{E}}^2 (1-\alpha_{\mathrm{G}})^4 k^2 + \mathcal{O}(k^4d_{\mathrm{nn}}^2).
\end{equation}

The two translational polarizations are degenerate at leading long-wavelength order because the homogenized FCC shear tensor is isotropic under
\begin{equation}
    \kappa_{\mathrm{n}} = \kappa_{\mathrm{t}}.
\end{equation}

Any splitting between the two transverse branches must arise from finite-wave-number cubic anisotropy, chiral polarization, nonuniform loading, higher-gradient contact terms, or additional phase--rotation coupling beyond the baseline parity-even constitutive order.

\section{Localized Defect--Condensate Relative-Flow Kinematics}

The coherent branch satisfies
\begin{equation}
    \dot{\mathbf{w}} = \dot{\mathbf{u}},
\end{equation}
and therefore
\begin{equation}
    \mathbf{v}_{\mathrm{rel}} = \dot{\mathbf{w}}-\dot{\mathbf{u}} = \mathbf{0}.
\end{equation}

A localized topological defect introduces a different kinematic sector. Translation, boundary deformation, or internal reconfiguration of the defect can induce condensate motion relative to the local FCC scaffold,
\begin{equation}
    \mathbf{v}_{\mathrm{rel}} \neq \mathbf{0}.
\end{equation}

Let $\Omega_{\mathrm{IC}}^{(\mathrm{def})}$ denote the connected interstitial-condensate domain influenced by the defect, and let $\partial\Omega_{\mathrm{def}}$ denote the effective defect--condensate interface.

A single localized defect mode is described by a generalized coordinate $q(t)$, restricted in the present construction to have dimensions of length,
\begin{equation}
    \boxed{ [q]=\mathrm{m} }.
\end{equation}
The restriction ensures that the boundary mode introduced below is dimensionless and that the associated surface-weighted measure has dimensions of area. Internal phase evolution or other dimensionless generalized coordinates require a separate coordinate metric and are not included in the present scalar translational normalization.

Let a material point on the defect boundary be represented by
\begin{equation}
    \mathbf{X}_{\partial} = \mathbf{X}_{\partial}(\boldsymbol{s},q),
\end{equation}
where $\boldsymbol{s}$ labels the boundary point. The boundary mode is
\begin{equation}
    \mathbf{b}_{\partial} \equiv \frac{\partial\mathbf{X}_{\partial}}{\partial q}.
\end{equation}

Since both $\mathbf{X}_{\partial}$ and $q$ have dimensions of length,
\begin{equation}
    [\mathbf{b}_{\partial}]=1.
\end{equation}

The defect-boundary velocity relative to the local scaffold is
\begin{equation}
    \dot{\mathbf{X}}_{\partial}^{(\mathrm{rel})} = \dot{q}\mathbf{b}_{\partial}.
\end{equation}

For a linearized localized response, write the condensate velocity relative to the scaffold as
\begin{equation}
    \mathbf{v}_{\mathrm{rel}}(\mathbf{x},t) = \dot{q}(t) \boldsymbol{\Upsilon}(\mathbf{x}),
\end{equation}
where $\boldsymbol{\Upsilon}(\mathbf{x})$ is the dimensionless condensate response field generated by unit generalized defect velocity.

The field $\boldsymbol{\Upsilon}$ depends on the defect geometry, the interstitial connectivity, the outer boundary conditions, and the selected defect mode. It is not assumed to be spatially uniform.

At conservative inviscid order, the defect boundary and the condensate satisfy the no-penetration condition
\begin{equation}
    \boxed{ \boldsymbol{\Upsilon}\cdot\hat{\mathbf{n}} = \mathbf{b}_{\partial}\cdot\hat{\mathbf{n}} \qquad \text{on } \partial\Omega_{\mathrm{def}} },
\end{equation}
where $\hat{\mathbf{n}}$ is directed from the defect into the condensate domain.

The no-penetration condition transmits only the normal boundary velocity. Tangential defect motion does not drive the conservative inviscid condensate response unless an additional tangential traction, interfacial locking, or no-slip law is introduced. The localized impedance projection must therefore be weighted by the normal boundary mode rather than by the full boundary-vector norm.

The coherent and localized sectors are consequently distinguished by
\begin{equation}
    \mathbf{v}_{\mathrm{rel}} = \mathbf{0} \quad \text{for the entrained periodic branch,}
\end{equation}
and
\begin{equation}
    \mathbf{v}_{\mathrm{rel}} = \dot{q}\boldsymbol{\Upsilon} \quad \text{for the localized defect-induced branch.}
\end{equation}

The kinetic contribution generated by $\boldsymbol{\Upsilon}$ is an added inertia of the defect coordinate. It is not an additional independent contribution to the coherent acoustic inertia.

\section{Conservative Boundary-Value Problem for Defect-Induced Condensate Flow}
\label{sec:AP_defect_boundary_value_problem}

The interstitial condensate is incrementally incompressible within the active FCC phase. The localized response field therefore satisfies
\begin{equation}
    \nabla\cdot\boldsymbol\Upsilon = 0
\end{equation}
throughout the defect-influenced condensate domain
\begin{equation}
    \Omega_{\mathrm{IC}}^{(\mathrm{def})}.
\end{equation}

On the minimum-energy zero-circulation branch, the localized condensate response is also irrotational,
\begin{equation}
    \nabla\times\boldsymbol\Upsilon = \mathbf 0.
\end{equation}
The response may then be represented by a reduced relative-flow potential
\begin{equation}
    \boldsymbol\Upsilon
    =
    \nabla\Phi_{\mathrm{rel}}.
\end{equation}
Since the response field is dimensionless,
\begin{equation}
    [\boldsymbol\Upsilon]=1,
\end{equation}
and
\begin{equation}
    [\nabla]=\mathrm m^{-1},
\end{equation}
the reduced potential has dimensions of length:
\begin{equation}
    [\Phi_{\mathrm{rel}}]=\mathrm m.
\end{equation}
The physical velocity potential associated with the defect motion is therefore
\begin{equation}
    \dot q\,\Phi_{\mathrm{rel}},
\end{equation}
which has the standard dimensions
\begin{equation}
    [\dot q\,\Phi_{\mathrm{rel}}]
    =
    \mathrm{m^2\,s^{-1}}.
\end{equation}

The incompressibility condition gives the Laplace equation
\begin{equation}
    \nabla^2\Phi_{\mathrm{rel}}=0
\end{equation}
within
\begin{equation}
    \Omega_{\mathrm{IC}}^{(\mathrm{def})}.
\end{equation}

At the effective defect--condensate interface, the conservative inviscid condition transmits only the normal boundary velocity. Therefore,
\begin{equation}
    \frac{\partial\Phi_{\mathrm{rel}}}{\partial n}
    =
    \mathbf b_{\partial}\cdot\hat{\mathbf n}
    \qquad
    \text{on }
    \partial\Omega_{\mathrm{def}}.
\end{equation}
Here \(\hat{\mathbf n}\) is directed from the defect into the condensate domain.

For an isolated defect in an unbounded condensate domain, the relative-flow field decays at large distance:
\begin{equation}
    \nabla\Phi_{\mathrm{rel}}
    \rightarrow
    \mathbf 0
    \qquad
    \text{as }
    |\mathbf x|\rightarrow\infty.
\end{equation}
For a periodic FCC defect calculation, \(\Phi_{\mathrm{rel}}\) satisfies the corresponding periodic boundary conditions together with zero net relative flux through the primitive-cell boundary.

Let the complete boundary of the defect-influenced condensate domain be
\begin{equation}
    \partial\Omega_{\mathrm{IC}}^{(\mathrm{def})}
    =
    \partial\Omega_{\mathrm{def}}
    \cup
    \partial\Omega_{\mathrm{out}}.
\end{equation}
Incompressibility requires the compatibility condition
\begin{equation}
    \int_{\partial\Omega_{\mathrm{IC}}^{(\mathrm{def})}}
    \boldsymbol\Upsilon\cdot\hat{\mathbf n}\,
    \mathrm dA
    =
    0.
\end{equation}
Equivalently,
\begin{equation}
    \int_{\partial\Omega_{\mathrm{def}}}
    \mathbf b_{\partial}\cdot\hat{\mathbf n}\,
    \mathrm dA
    +
    \int_{\partial\Omega_{\mathrm{out}}}
    \boldsymbol\Upsilon\cdot\hat{\mathbf n}\,
    \mathrm dA
    =
    0.
\end{equation}

For an isolated defect with vanishing far-field flux, the admissible boundary mode must satisfy
\begin{equation}
    \int_{\partial\Omega_{\mathrm{def}}}
    \mathbf b_{\partial}\cdot\hat{\mathbf n}\,
    \mathrm dA
    =
    0.
\end{equation}
A rigid translation satisfies this condition because
\begin{equation}
    \int_{\partial\Omega_{\mathrm{def}}}
    \hat{\mathbf n}\,\mathrm dA
    =
    \mathbf 0.
\end{equation}
A purely radial breathing mode that changes the enclosed defect volume does not satisfy the isolated-domain compatibility condition unless its flux is balanced by an outer-boundary flow, a compensating scaffold deformation, or a coupled phase-conversion process. Such a mode is therefore not an admissible isolated perturbation of the strictly incompressible conservative branch.

The conservative response field may equivalently be obtained from the constrained minimization problem
\begin{equation}
    \boldsymbol\Upsilon
    =
    \operatorname*{arg\,min}_{\mathbf v}
    \int_{\Omega_{\mathrm{IC}}^{(\mathrm{def})}}
    |\mathbf v|^2\,\mathrm dV,
\end{equation}
subject to
\begin{equation}
    \nabla\cdot\mathbf v=0,
\end{equation}
and
\begin{equation}
    \mathbf v\cdot\hat{\mathbf n}
    =
    \mathbf b_{\partial}\cdot\hat{\mathbf n}
    \qquad
    \text{on }
    \partial\Omega_{\mathrm{def}},
\end{equation}
together with the selected outer-boundary condition.

The zero-circulation restriction selects the minimum-energy potential-flow branch. A defect carrying conserved circulation, vortex content, or chiral winding requires an additional solenoidal harmonic component,
\begin{equation}
    \boldsymbol\Upsilon
    =
    \nabla\Phi_{\mathrm{rel}}
    +
    \boldsymbol\Upsilon_{\mathrm{circ}},
\end{equation}
with
\begin{equation}
    \nabla\cdot\boldsymbol\Upsilon_{\mathrm{circ}}=0,
    \qquad
    \nabla\times\boldsymbol\Upsilon_{\mathrm{circ}}\neq\mathbf 0.
\end{equation}
Such circulation-carrying branches are not included in the baseline impedance projection developed here.

For a prescribed admissible defect geometry and boundary mode, the boundary-value problem determines \(\boldsymbol\Upsilon\) uniquely up to the circulation sector and the irrelevant additive constant in \(\Phi_{\mathrm{rel}}\). The resulting field supplies the mode-weighted condensate participation volume and the localized added-inertia coefficient developed below.

\section{Normal Boundary Weighting and Microscopic Impedance Projection}
\label{sec:AP_normal_boundary_impedance}

The conservative boundary-value problem transmits the normal component of the defect-boundary velocity into the interstitial condensate. The appropriate interface measure must therefore be constructed from the normal boundary mode,
\begin{equation}
    b_n(\boldsymbol{s})
    \equiv
    \mathbf{b}_{\partial}(\boldsymbol{s})
    \cdot
    \hat{\mathbf{n}}(\boldsymbol{s}).
\end{equation}

The mode-weighted normal interface area is defined by
\begin{equation}
    A_{\mathrm{imp}}^{(n)}
    \equiv
    \int_{\partial\Omega_{\mathrm{def}}}
    \left|
    \mathbf{b}_{\partial}\cdot\hat{\mathbf{n}}
    \right|^2
    \mathrm{d}A .
\end{equation}

Since \(\mathbf{b}_{\partial}\) is dimensionless,
\begin{equation}
    \left[
    A_{\mathrm{imp}}^{(n)}
    \right]
    =
    \mathrm{m}^2 .
\end{equation}

The quantity \(A_{\mathrm{imp}}^{(n)}\) is not, in general, equal to the total geometric area of the defect boundary. It is a mode-dependent measure of the portion of the interface that participates in the normal displacement transmitted to the condensate. Boundary regions with vanishing normal motion do not contribute. Regions with larger normal displacement relative to the generalized coordinate \(q\) receive correspondingly greater weight.

The condensate participation volume is defined by
\begin{equation}
    V_{\mathrm{imp}}
    \equiv
    \int_{\Omega_{\mathrm{IC}}^{(\mathrm{def})}}
    \left|
    \boldsymbol{\Upsilon}
    \right|^2
    \mathrm{d}V .
\end{equation}

Because \(\boldsymbol{\Upsilon}\) is dimensionless,
\begin{equation}
    \left[
    V_{\mathrm{imp}}
    \right]
    =
    \mathrm{m}^3 .
\end{equation}

The quantity \(V_{\mathrm{imp}}\) is an inertia-weighted participation volume rather than the geometric volume of the entire connected interstitial domain. Regions in which the relative-flow amplitude is negligible contribute negligibly to the kinetic functional even when they are geometrically connected to the defect. Conversely, a weak but spatially extended response may contribute substantially after integration over the condensate domain.

The localized surface-to-volume impedance projection is defined by
\begin{equation}
    \label{eq:AP_defect_projection}
    \iota_V^{(\mathrm{def})}
    \equiv
    \frac{
    A_{\mathrm{imp}}^{(n)}
    }{
    V_{\mathrm{imp}}
    } .
\end{equation}

Its dimensions are
\begin{equation}
    \left[
    \iota_V^{(\mathrm{def})}
    \right]
    =
    \mathrm{m}^{-1}.
\end{equation}

This localized defect projection must be distinguished from the coherent transverse modal projection of the periodic FCC transfer mode. The coherent projection is
\begin{equation}
    \iota_{\mathrm{tr}}^{(\mathrm{modal})}
    =
    \frac{
    A_{\mathrm{ct}}^{(\mathrm{modal})}
    }{
    V_p
    }
    =
    \frac{
    \rho_{\mathrm{eff}}^{(\mathrm{micro})}
    }{
    \mu_{\mathrm D}
    } .
\end{equation}

The localized defect-flow projection is instead
\begin{equation}
    \iota_V^{(\mathrm{def})}
    =
    \frac{
    A_{\mathrm{imp}}^{(n)}
    }{
    V_{\mathrm{imp}}
    } .
\end{equation}

Therefore,
\begin{equation}
    \label{eq:AP_projection_distinction}
    \iota_V^{(\mathrm{def})}
    \neq
    \iota_{\mathrm{tr}}^{(\mathrm{modal})}
\end{equation}
unless the periodic coherent-mode integral and the localized defect-flow integral are independently shown to coincide.

The equality of dimensions is not a mechanical identification. The quantity \(\iota_{\mathrm{tr}}^{(\mathrm{modal})}\) belongs to the entrained coherent transverse acoustic sector, whereas \(\iota_V^{(\mathrm{def})}\) belongs to a boundary-value problem for relative condensate flow around a localized defect.

In terms of the solved boundary and volume fields,
\begin{equation}
    \iota_V^{(\mathrm{def})}
    =
    \frac{
    \displaystyle
    \int_{\partial\Omega_{\mathrm{def}}}
    \left|
    \mathbf{b}_{\partial}\cdot\hat{\mathbf{n}}
    \right|^2
    \mathrm{d}A
    }{
    \displaystyle
    \int_{\Omega_{\mathrm{IC}}^{(\mathrm{def})}}
    \left|
    \boldsymbol{\Upsilon}
    \right|^2
    \mathrm{d}V
    } .
\end{equation}

This expression gives the projection a definite mechanical meaning. It is the mode-weighted inverse length that converts the surface inertia of the active interface into the volumetric impedance sampled by the localized condensate response.

The normalization of the generalized coordinate does not affect the projection. Under the rescaling
\begin{equation}
    q^{\prime}=a q,
\end{equation}
the boundary mode transforms as
\begin{equation}
    \mathbf{b}_{\partial}^{\prime}
    =
    \frac{1}{a}
    \mathbf{b}_{\partial},
\end{equation}
and the condensate response field transforms as
\begin{equation}
    \boldsymbol{\Upsilon}^{\prime}
    =
    \frac{1}{a}
    \boldsymbol{\Upsilon}.
\end{equation}

Consequently,
\begin{equation}
    A_{\mathrm{imp}}^{(n)\prime}
    =
    \frac{1}{a^2}
    A_{\mathrm{imp}}^{(n)},
    \qquad
    V_{\mathrm{imp}}^{\prime}
    =
    \frac{1}{a^2}
    V_{\mathrm{imp}} .
\end{equation}

Therefore,
\begin{equation}
    \iota_V^{(\mathrm{def})\prime}
    =
    \iota_V^{(\mathrm{def})}.
\end{equation}

The projection is thus a property of the physical defect mode and its induced condensate field, not of the arbitrary normalization assigned to \(q\).

Define the associated impedance penetration length by
\begin{equation}
    \ell_{\mathrm{imp}}^{(\mathrm{def})}
    \equiv
    \frac{
    V_{\mathrm{imp}}
    }{
    A_{\mathrm{imp}}^{(n)}
    } .
\end{equation}

Then
\begin{equation}
    \iota_V^{(\mathrm{def})}
    =
    \frac{
    1
    }{
    \ell_{\mathrm{imp}}^{(\mathrm{def})}
    } .
\end{equation}

The penetration length characterizes the effective depth over which condensate motion contributes to the added inertia of the selected defect mode. It is not necessarily identical to the geometric radius, shell thickness, or lattice spacing of the defect. Its value is determined by the complete solution of the conservative boundary-value problem.

The earlier empirical convention
\begin{equation}
    \iota_V
    =
    1\ \mathrm{m}^{-1}
\end{equation}
corresponds to
\begin{equation}
    \ell_{\mathrm{imp}}
    =
    1\ \mathrm{m}.
\end{equation}

Within the present formulation, that assignment remains an empirical calibration until a specific defect geometry and response field yield the same value through the ratio of the mode integrals.

\section{Mode-Resolved Condensate Impedance Density}

The active-vacuum surface inertia is denoted by $\mu_{\mathrm{D}}$. For a localized defect mode with generalized velocity $\dot{q}$, the interface representation of the conservative kinetic energy is
\begin{equation}
    T_{\partial} = \frac{1}{2} \mu_{\mathrm{D}} A_{\mathrm{imp}}^{(n)} \dot{q}^2,
\end{equation}
before any additional phase-dependent activation is imposed.

The same kinetic energy may be represented through an effective volumetric condensate impedance density,
\begin{equation}
    T_{\mathrm{vol}} = \frac{1}{2} \rho_{\mathrm{IC}}^{(\mathrm{imp})} V_{\mathrm{imp}} \dot{q}^2.
\end{equation}

Equality of the two representations requires
\begin{equation}
    \mu_{\mathrm{D}} A_{\mathrm{imp}}^{(n)} = \rho_{\mathrm{IC}}^{(\mathrm{imp})} V_{\mathrm{imp}}.
\end{equation}

Therefore,
\begin{equation}
    \rho_{\mathrm{IC}}^{(\mathrm{imp})} = \mu_{\mathrm{D}} \frac{ A_{\mathrm{imp}}^{(n)} }{ V_{\mathrm{imp}} }.
\end{equation}

Using the microscopic projection,
\begin{equation}
    \boxed{ \rho_{\mathrm{IC}}^{(\mathrm{imp})} = \mu_{\mathrm{D}} \iota_V^{(\mathrm{def})} }.
\end{equation}

Equivalently,
\begin{equation}
    \rho_{\mathrm{IC}}^{(\mathrm{imp})} = \frac{ \mu_{\mathrm{D}} }{ \ell_{\mathrm{imp}}^{(\mathrm{def})} }.
\end{equation}

The dimensions are $[\rho_{\mathrm{IC}}^{(\mathrm{imp})}] = \mathrm{kg\,m^{-3}}$ because $[\mu_{\mathrm{D}}] = \mathrm{kg\,m^{-2}}$ and $[\iota_V^{(\mathrm{def})}] = \mathrm{m}^{-1}$.

The quantity $\rho_{\mathrm{IC}}^{(\mathrm{imp})}$ is a mode-resolved impedance density. It is not automatically the intrinsic equilibrium density of the interstitial condensate. It measures the volumetric inertia sampled by the specific localized flow field $\boldsymbol{\Upsilon}$.

The associated added-inertia coefficient is
\begin{equation}
    m_{\mathrm{add}} = \rho_{\mathrm{IC}}^{(\mathrm{imp})} V_{\mathrm{imp}}.
\end{equation}

Using the surface representation, $m_{\mathrm{add}} = \mu_{\mathrm{D}} A_{\mathrm{imp}}^{(n)}$. Thus,
\begin{equation}
    m_{\mathrm{add}} = \mu_{\mathrm{D}} A_{\mathrm{imp}}^{(n)} = \rho_{\mathrm{IC}}^{(\mathrm{imp})} V_{\mathrm{imp}}.
\end{equation}

The added inertia is the quadratic kinetic coefficient of the selected defect coordinate. It is generated by the condensate flow required to satisfy the boundary motion and incompressibility constraints. It is not the static material mass of the defect core.

The distinction between coherent and localized inertia is therefore $\rho_{\mathrm{eff}}^{(0)}$ for the entrained periodic coherent coordinate, and $\rho_{\mathrm{IC}}^{(\mathrm{imp})}$ for the localized relative-flow coordinate.

The coherent inertia is assigned to the common scaffold--condensate displacement field, $\mathbf{v}_{\mathrm{rel}}=\mathbf{0}$. The localized impedance density is assigned to the relative-flow field, $\mathbf{v}_{\mathrm{rel}} = \dot{q}\boldsymbol{\Upsilon}$. These coefficients belong to different kinetic reductions and must not be added to the same generalized coordinate without an explicit coupled two-field action.

The projection may depend on the defect class and on the selected mode:
\begin{equation}
    \iota_V^{(\mathrm{def})} = \iota_V^{(\mathrm{def})} \left( \text{geometry}, \text{topology}, \text{boundary mode}, \text{outer boundary condition} \right).
\end{equation}

At finite frequency, the condensate response may also become dispersive. A frequency-dependent projection may then be written as $\iota_V^{(\mathrm{def})}(\omega)$. The conservative low-frequency impedance is defined by
\begin{equation}
    \iota_{V,0}^{(\mathrm{def})} \equiv \lim_{\omega\rightarrow0} \operatorname{Re} \left[ \iota_V^{(\mathrm{def})}(\omega) \right].
\end{equation}

When dissipative or radiative effects are included, the corresponding dynamic impedance may acquire an imaginary component. The imaginary part represents loss or phase lag and lies outside the conservative action developed here.

A universal condensate impedance density may be introduced only if the low-frequency projection is common to an identified universality class of stable defects:
\begin{equation}
    \iota_{V,0}^{(\mathrm{def})} = \iota_{V,0}
\end{equation}
for all defects in that class. Under that condition, $\rho_{\mathrm{IC}} \equiv \mu_{\mathrm{D}}\iota_{V,0}$.

If the projection depends materially on defect topology, boundary geometry, or mode structure, then the notation $\rho_{\mathrm{IC}}^{(\mathrm{imp})}$ must be retained, and the quantity must not be interpreted as a universal intrinsic condensate density.

The microscopic projection derived here therefore replaces an unspecified dimensional conversion by a boundary-value functional. It does not, by itself, establish universality or assign a numerical value.

\section{Reduced Swept-Volume Approximation for Localized Defects}

The exact localized added inertia is determined by the mode-weighted volume
\begin{equation}
    V_{\mathrm{imp}} = \int_{\Omega_{\mathrm{IC}}^{(\mathrm{def})}} |\boldsymbol{\Upsilon}|^2 \mathrm{d}V.
\end{equation}

For a thin localized defect layer surrounding a spherical Soft Floor node, this volume may be related to a geometric swept-volume approximation. Let the reference node radius be $r_{\mathrm{D}}$, and let $r_{\mathrm{amp}}$ denote the equilibrium radial thickness associated with the localized defect structure. The geometric shell volume is
\begin{equation}
    V_{\mathrm{sweep}} = \frac{4\pi}{3} \left[ (r_{\mathrm{D}}+r_{\mathrm{amp}})^3 - r_{\mathrm{D}}^3 \right].
\end{equation}

Expansion gives
\begin{equation}
    V_{\mathrm{sweep}} = 4\pi r_{\mathrm{D}}^2r_{\mathrm{amp}} + 4\pi r_{\mathrm{D}}r_{\mathrm{amp}}^2 + \frac{4\pi}{3}r_{\mathrm{amp}}^3.
\end{equation}

In the thin-layer regime ($r_{\mathrm{amp}} \ll r_{\mathrm{D}}$), the leading term is
\begin{equation}
    V_{\mathrm{sweep}} = 4\pi r_{\mathrm{D}}^2r_{\mathrm{amp}} + \mathcal{O} \left( r_{\mathrm{D}}r_{\mathrm{amp}}^2 \right).
\end{equation}

The familiar expression
\begin{equation}
    V_{\mathrm{eff}} = 4\pi r_{\mathrm{D}}^2r_{\mathrm{amp}}
\end{equation}
is therefore the leading geometric shell-volume approximation.

The inertia-weighted participation volume $V_{\mathrm{imp}}$ need not coincide with $V_{\mathrm{sweep}}$. The relation between them is represented by the dimensionless geometry factor
\begin{equation}
    \mathcal{G}_{\mathrm{imp}} \equiv \frac{ V_{\mathrm{imp}} }{ V_{\mathrm{sweep}} }.
\end{equation}

Thus,
\begin{equation}
    V_{\mathrm{imp}} = \mathcal{G}_{\mathrm{imp}} V_{\mathrm{sweep}}.
\end{equation}

In the thin-layer limit,
\begin{equation}
    V_{\mathrm{imp}} = 4\pi \mathcal{G}_{\mathrm{imp}} r_{\mathrm{D}}^2 r_{\mathrm{amp}} + \mathcal{O} \left( r_{\mathrm{D}}r_{\mathrm{amp}}^2 \right).
\end{equation}

The factor $\mathcal{G}_{\mathrm{imp}}$ accounts for the difference between the geometric shell volume and the kinetic participation volume determined by the solved condensate field. It may exceed, equal, or fall below unity, depending on the spatial extent and amplitude distribution of $\boldsymbol{\Upsilon}$.

The condition $\mathcal{G}_{\mathrm{imp}}=1$ is not a geometric identity. It is a constitutive property of the induced condensate response and must be demonstrated for the relevant defect mode.

Substitution into the added-inertia relation gives $m_{\mathrm{add}} = \rho_{\mathrm{IC}}^{(\mathrm{imp})} V_{\mathrm{imp}}$. Therefore,
\begin{equation}
    m_{\mathrm{add}} = 4\pi r_{\mathrm{D}}^2 \rho_{\mathrm{IC}}^{(\mathrm{imp})} \mathcal{G}_{\mathrm{imp}} r_{\mathrm{amp}} + \mathcal{O} \left( r_{\mathrm{D}}r_{\mathrm{amp}}^2 \right).
\end{equation}

Define the mode-resolved defect impedance coefficient by
\begin{equation}
    \boxed{ K_{\mathrm{D}}^{(\mathrm{def})} \equiv 4\pi r_{\mathrm{D}}^2 \rho_{\mathrm{IC}}^{(\mathrm{imp})} \mathcal{G}_{\mathrm{imp}} }.
\end{equation}

Its dimensions are
\begin{equation}
    \left[ K_{\mathrm{D}}^{(\mathrm{def})} \right] = \mathrm{kg\,m^{-1}}.
\end{equation}

The leading reduced added-inertia law is then
\begin{equation}
    \boxed{ m_{\mathrm{add}} = K_{\mathrm{D}}^{(\mathrm{def})} r_{\mathrm{amp}} + \mathcal{O} \left( r_{\mathrm{amp}}^2 \right) }.
\end{equation}

Using $\rho_{\mathrm{IC}}^{(\mathrm{imp})} = \mu_{\mathrm{D}} \iota_V^{(\mathrm{def})}$, the impedance coefficient becomes
\begin{equation}
    K_{\mathrm{D}}^{(\mathrm{def})} = 4\pi r_{\mathrm{D}}^2 \mu_{\mathrm{D}} \iota_V^{(\mathrm{def})} \mathcal{G}_{\mathrm{imp}}.
\end{equation}

The reduced law therefore follows from three independent ingredients:
\begin{itemize}
    \item $4\pi r_{\mathrm{D}}^2r_{\mathrm{amp}}$ as the leading geometric swept volume,
    \item $\iota_V^{(\mathrm{def})}$ as the mode-resolved surface-to-volume impedance projection, and
    \item $\mathcal{G}_{\mathrm{imp}}$ as the correction from geometric swept volume to kinetic participation volume.
\end{itemize}

The existing impedance relation
\begin{equation}
    K_{\mathrm{D}} = 4\pi r_{\mathrm{D}}^2\rho_{\mathrm{IC}}
\end{equation}
is recovered only when $\rho_{\mathrm{IC}}^{(\mathrm{imp})} = \rho_{\mathrm{IC}}$ is universal within the selected defect class and $\mathcal{G}_{\mathrm{imp}}=1$. The geometric swept-volume approximation does not independently establish either condition.

\section{Phase-Participation Function and Branch-Level Sine Closure}

The localized condensate response determined in the previous sections fixes the hydrodynamic inertia available to a prescribed defect mode. The fraction of that inertia activated by the defect phase structure is represented by the dimensionless participation function
\begin{equation}
    f_{\varphi}(\Delta\varphi).
\end{equation}

When normalized to the fully activated state, the function satisfies
\begin{equation}
    0 \leq f_{\varphi}(\Delta\varphi) \leq 1.
\end{equation}

The nonnegativity condition follows from its occurrence as a coefficient of a quadratic kinetic term. A negative value would render the localized kinetic Hessian indefinite and is therefore inadmissible.

For a spatially uniform phase mismatch across the active defect boundary, the phase-weighted kinetic energy is
\begin{equation}
    T_{\mathrm{def}} = \frac{1}{2} \mu_{\mathrm{D}} A_{\mathrm{imp}}^{(n)} f_{\varphi}(\Delta\varphi) \dot{q}^2.
\end{equation}

Using
\begin{equation}
    \rho_{\mathrm{IC}}^{(\mathrm{imp})} V_{\mathrm{imp}} = \mu_{\mathrm{D}} A_{\mathrm{imp}}^{(n)},
\end{equation}
the equivalent volumetric representation is
\begin{equation}
    T_{\mathrm{def}} = \frac{1}{2} \rho_{\mathrm{IC}}^{(\mathrm{imp})} V_{\mathrm{imp}} f_{\varphi}(\Delta\varphi) \dot{q}^2.
\end{equation}

The corresponding phase-activated topological loading is
\begin{equation}
    M_{\mathrm{top}}(\Delta\varphi) = \mu_{\mathrm{D}} A_{\mathrm{imp}}^{(n)} f_{\varphi}(\Delta\varphi),
\end{equation}
or, equivalently,
\begin{equation}
    M_{\mathrm{top}}(\Delta\varphi) = \rho_{\mathrm{IC}}^{(\mathrm{imp})} V_{\mathrm{imp}} f_{\varphi}(\Delta\varphi).
\end{equation}

The phase-participation function is periodic,
\begin{equation}
    f_{\varphi}(\Delta\varphi+2\pi) = f_{\varphi}(\Delta\varphi).
\end{equation}

Within an unpolarized reciprocal defect sector, reversal of the phase mismatch does not alter the kinetic loading, and therefore
\begin{equation}
    f_{\varphi}(-\Delta\varphi) = f_{\varphi}(\Delta\varphi).
\end{equation}

This evenness condition is not imposed on an individual chiral branch. A polarized defect may distinguish the two signs of $\Delta\varphi$, provided its kinetic coefficient remains nonnegative on the selected branch.

The condition
\begin{equation}
    f_{\varphi}(0) = 0
\end{equation}
is valid only when phase matching eliminates the localized defect boundary and therefore removes the associated relative-flow obstruction. In that defect class,
\begin{equation}
    \Delta\varphi = 0
\end{equation}
implies
\begin{equation}
    A_{\mathrm{imp}}^{(n)} = 0, \qquad V_{\mathrm{imp}} = 0,
\end{equation}
in the physical matched-phase limit, and hence
\begin{equation}
    M_{\mathrm{top}}(0) = 0.
\end{equation}

A boundary that remains geometrically present at $\Delta\varphi=0$ possesses a nonzero phase-independent added inertia. Its total kinetic coefficient must then be decomposed as
\begin{equation}
    M_{\mathrm{add}}(\Delta\varphi) = M_{\mathrm{geom}} + M_{\varphi} f_{\varphi}(\Delta\varphi),
\end{equation}
where
\begin{equation}
    M_{\mathrm{geom}} > 0
\end{equation}
is the phase-independent geometric contribution and
\begin{equation}
    f_{\varphi}(0) = 0
\end{equation}
applies only to the phase-activated increment. The purely multiplicative law
\begin{equation}
    M_{\mathrm{top}} = M_{\mathrm{add}}^{(0)} f_{\varphi}
\end{equation}
is therefore restricted to phase-created defect boundaries.

The reduced sine closure used in the topological mass law is admissible only as a branch-level representation. Introduce the branch orientation
\begin{equation}
    \chi = \pm 1.
\end{equation}

The nonnegative sine closure is
\begin{equation}
    f_{\varphi,\chi}^{(\sin)}(\Delta\varphi) = \chi\sin(\Delta\varphi),
\end{equation}
restricted to the branch domain
\begin{equation}
    \chi\sin(\Delta\varphi) \geq 0.
\end{equation}

For the positive branch,
\begin{equation}
    0 \leq \Delta\varphi \leq \pi,
\end{equation}
one has
\begin{equation}
    f_{\varphi,+}^{(\sin)} = \sin(\Delta\varphi).
\end{equation}

For the reciprocal negative branch,
\begin{equation}
    -\pi \leq \Delta\varphi \leq 0,
\end{equation}
one has
\begin{equation}
    f_{\varphi,-}^{(\sin)} = -\sin(\Delta\varphi).
\end{equation}

Under simultaneous reversal,
\begin{equation}
    \chi \rightarrow -\chi, \qquad \Delta\varphi \rightarrow -\Delta\varphi,
\end{equation}
the participation factor remains invariant:
\begin{equation}
    f_{\varphi,-\chi}^{(\sin)}(-\Delta\varphi) = f_{\varphi,\chi}^{(\sin)}(\Delta\varphi).
\end{equation}

The reciprocal branch envelope is therefore
\begin{equation}
    f_{\varphi}^{(\mathrm{rec})}(\Delta\varphi) = |\sin(\Delta\varphi)|.
\end{equation}

The unsigned envelope is appropriate for a kinetic coefficient because the added inertia cannot change sign. Any orientation carried by the sign of $\sin(\Delta\varphi)$ belongs to a chiral force, circulation, or topological-orientation sector rather than to the positive quadratic mass coefficient.

The sine closure is consequently a reduced constitutive approximation on a specified phase branch. Its global extension, branch junctions, saturation behavior, and relation to the nonlinear phase-contact potential must be supplied by the complete defect phase action.

When the phase mismatch varies over the defect boundary, the participation factor cannot be extracted from the surface integral. The kinetic energy must instead be written as
\begin{equation}
    T_{\mathrm{def}} = \frac{1}{2} \mu_{\mathrm{D}} \dot{q}^2 \int_{\partial\Omega_{\mathrm{def}}} \left| \mathbf{b}_{\partial}\cdot\hat{\mathbf{n}} \right|^2 f_{\varphi} \!\left[ \Delta\varphi(\boldsymbol{s}) \right] \mathrm{d}A.
\end{equation}

The scalar reduction
\begin{equation}
    T_{\mathrm{def}} = \frac{1}{2} \mu_{\mathrm{D}} A_{\mathrm{imp}}^{(n)} f_{\varphi}(\Delta\varphi) \dot{q}^2
\end{equation}
therefore assumes a phase mismatch that is uniform, or effectively uniform under the normal-mode weighting, over the active defect interface.

\section{Recovery Conditions for the Hydrodynamic Mass Law}

The mode-resolved topological loading is
\begin{equation}
    M_{\mathrm{top}} = \rho_{\mathrm{IC}}^{(\mathrm{imp})} V_{\mathrm{imp}} f_{\varphi}(\Delta\varphi).
\end{equation}

Using the relation
\begin{equation}
    V_{\mathrm{imp}} = \mathcal{G}_{\mathrm{imp}} V_{\mathrm{sweep}},
\end{equation}
and the thin-layer expansion
\begin{equation}
    V_{\mathrm{sweep}} = 4\pi r_{\mathrm{D}}^2r_{\mathrm{amp}} + \mathcal{O} \left( r_{\mathrm{D}}r_{\mathrm{amp}}^2 \right),
\end{equation}
one obtains
\begin{equation}
    M_{\mathrm{top}} = 4\pi r_{\mathrm{D}}^2 \rho_{\mathrm{IC}}^{(\mathrm{imp})} \mathcal{G}_{\mathrm{imp}} r_{\mathrm{amp}} f_{\varphi}(\Delta\varphi) + \mathcal{O} \left( r_{\mathrm{D}}r_{\mathrm{amp}}^2 \right).
\end{equation}

With the mode-resolved impedance coefficient
\begin{equation}
    K_{\mathrm{D}}^{(\mathrm{def})} = 4\pi r_{\mathrm{D}}^2 \rho_{\mathrm{IC}}^{(\mathrm{imp})} \mathcal{G}_{\mathrm{imp}},
\end{equation}
the reduced mass law becomes
\begin{equation}
    M_{\mathrm{top}} = K_{\mathrm{D}}^{(\mathrm{def})} r_{\mathrm{amp}} f_{\varphi}(\Delta\varphi) + \mathcal{O} \left( r_{\mathrm{amp}}^2 \right).
\end{equation}

Using the surface-to-volume projection,
\begin{equation}
    \rho_{\mathrm{IC}}^{(\mathrm{imp})} = \mu_{\mathrm{D}} \iota_V^{(\mathrm{def})},
\end{equation}
the impedance coefficient is
\begin{equation}
    K_{\mathrm{D}}^{(\mathrm{def})} = 4\pi r_{\mathrm{D}}^2 \mu_{\mathrm{D}} \iota_V^{(\mathrm{def})} \mathcal{G}_{\mathrm{imp}}.
\end{equation}

The previously employed hydrodynamic mass law is recovered under the simultaneous conditions
\begin{equation}
    r_{\mathrm{amp}} \ll r_{\mathrm{D}}, \qquad \mathcal{G}_{\mathrm{imp}} = 1, \qquad \rho_{\mathrm{IC}}^{(\mathrm{imp})} = \rho_{\mathrm{IC}},
\end{equation}
and
\begin{equation}
    f_{\varphi}(\Delta\varphi) = \sin(\Delta\varphi)
\end{equation}
on the selected positive phase branch
\begin{equation}
    0 \leq \Delta\varphi \leq \pi.
\end{equation}

Under these restrictions,
\begin{equation}
    V_{\mathrm{imp}} = V_{\mathrm{eff}} = 4\pi r_{\mathrm{D}}^2r_{\mathrm{amp}},
\end{equation}
and therefore
\begin{equation}
    M_{\mathrm{top}} = \rho_{\mathrm{IC}} V_{\mathrm{eff}} \sin(\Delta\varphi).
\end{equation}

Defining
\begin{equation}
    K_{\mathrm{D}} = 4\pi r_{\mathrm{D}}^2\rho_{\mathrm{IC}},
\end{equation}
one obtains
\begin{equation}
    M_{\mathrm{top}} = K_{\mathrm{D}} r_{\mathrm{amp}} \sin(\Delta\varphi).
\end{equation}

The formula
\begin{equation}
    M_{\mathrm{top}} = \rho_{\mathrm{IC}} V_{\mathrm{eff}} \sin(\Delta\varphi)
\end{equation}
is therefore not the unrestricted microscopic law. It is the thin-layer, uniform-phase, universal-projection, unit-geometry-factor reduction of the mode-resolved expression
\begin{equation}
    M_{\mathrm{top}} = \rho_{\mathrm{IC}}^{(\mathrm{imp})} V_{\mathrm{imp}} f_{\varphi}(\Delta\varphi).
\end{equation}

For the reciprocal negative phase branch, positivity requires
\begin{equation}
    M_{\mathrm{top}} = -\rho_{\mathrm{IC}} V_{\mathrm{eff}} \sin(\Delta\varphi),
\end{equation}
with
\begin{equation}
    -\pi \leq \Delta\varphi \leq 0.
\end{equation}

Equivalently, the reciprocal two-branch form is
\begin{equation}
    M_{\mathrm{top}} = \rho_{\mathrm{IC}} V_{\mathrm{eff}} |\sin(\Delta\varphi)|.
\end{equation}

The signed sine may continue to encode branch orientation in a chiral coupling, but the inertial loading entering the kinetic action remains nonnegative.

The reduced law also assumes that the defect boundary is phase generated. When a persistent impermeable boundary remains at phase matching, the total inertia is instead
\begin{equation}
    M_{\mathrm{add}} = M_{\mathrm{geom}} + \rho_{\mathrm{IC}}^{(\mathrm{imp})} V_{\mathrm{imp}} f_{\varphi}(\Delta\varphi).
\end{equation}

In that case, the hydrodynamic sine law describes only the phase-dependent increment and not the complete added inertia.

When the topological loading is used as the inertial coefficient of the localized defect coordinate,
\begin{equation}
    m_i = M_{\mathrm{top}},
\end{equation}
and as the coefficient coupling the same defect to the restorative tension potential,
\begin{equation}
    \mathbf{F}_{\mathrm{grav}} = -M_{\mathrm{top}} \nabla\Phi_{\tau},
\end{equation}
the acceleration is
\begin{equation}
    \mathbf{a} = \frac{\mathbf{F}_{\mathrm{grav}}}{m_i} = -\nabla\Phi_{\tau}.
\end{equation}

The cancellation requires the same mode-resolved loading functional in both operations. Replacing $M_{\mathrm{top}}$ by independently normalized inertial and restorative coefficients would remove the mechanical identity on which the cancellation depends.

The hydrodynamic mass law is thereby recovered as a controlled reduction of the defect–condensate boundary-value formulation rather than as an independent dimensional postulate. 

\section{Consequences for the Empirical Projection and Defect-Flow Calibration}

The localized surface-to-volume projection is the functional
\begin{equation}
    \iota_V^{(\mathrm{def})} = \frac{ \displaystyle \int_{\partial\Omega_{\mathrm{def}}} \left| \mathbf{b}_{\partial}\cdot\hat{\mathbf{n}} \right|^2 \mathrm{d}A }{ \displaystyle \int_{\Omega_{\mathrm{IC}}^{(\mathrm{def})}} |\boldsymbol{\Upsilon}|^2 \mathrm{d}V }.
\end{equation}

Its value depends on the complete defect-flow problem. More explicitly,
\begin{equation}
    \iota_V^{(\mathrm{def})} = \iota_V^{(\mathrm{def})} \left[ \partial\Omega_{\mathrm{def}}, \mathbf{b}_{\partial}, \Omega_{\mathrm{IC}}^{(\mathrm{def})}, \mathcal{B}_{\mathrm{out}}, \mathcal{T}_{\mathrm{def}} \right],
\end{equation}
where $\mathcal{B}_{\mathrm{out}}$ denotes the outer-boundary condition and $\mathcal{T}_{\mathrm{def}}$ denotes the topological and phase class of the defect.

The response field entering the denominator is determined by
\begin{equation}
    \nabla\cdot\boldsymbol{\Upsilon} = 0,
\end{equation}
together with
\begin{equation}
    \boldsymbol{\Upsilon}\cdot\hat{\mathbf{n}} = \mathbf{b}_{\partial}\cdot\hat{\mathbf{n}}
\end{equation}
on the defect boundary and the selected condition on the outer boundary. On the zero-circulation conservative branch,
\begin{equation}
    \boldsymbol\Upsilon=\nabla\Phi_{\mathrm{rel}},
\end{equation}
with
\begin{equation}
    \nabla^2\Phi_{\mathrm{rel}}=0.
\end{equation}

Consequently, the definition of $\iota_V^{(\mathrm{def})}$ is complete only as a boundary-value functional. It does not provide a numerical value until the corresponding response field has been solved for a specified defect geometry and mode.

The relation
\begin{equation}
    \rho_{\mathrm{IC}}^{(\mathrm{imp})} = \mu_{\mathrm{D}} \iota_V^{(\mathrm{def})}
\end{equation}
is therefore a derived conversion between the surface inertia and the mode-resolved volumetric impedance. The numerical assignment
\begin{equation}
    \iota_V = 1\ \mathrm{m}^{-1}
\end{equation}
is not derived by the present appendix. It remains an empirical normalization corresponding to the penetration length
\begin{equation}
    \ell_{\mathrm{imp}} = 1\ \mathrm{m}.
\end{equation}

Under this empirical assignment,
\begin{equation}
    \rho_{\mathrm{IC}}^{(\mathrm{emp})} = \mu_{\mathrm{D}} \left( 1\ \mathrm{m}^{-1} \right).
\end{equation}

For the reduced thin-layer defect law with
\begin{equation}
    \mathcal{G}_{\mathrm{imp}} = 1,
\end{equation}
the corresponding empirical impedance coefficient is
\begin{equation}
    K_{\mathrm{D}}^{(\mathrm{emp})} = 4\pi r_{\mathrm{D}}^2 \mu_{\mathrm{D}} \left( 1\ \mathrm{m}^{-1} \right).
\end{equation}

Any calibration of $r_{\mathrm{D}}$ that uses
\begin{equation}
    \rho_{\mathrm{IC}} = \mu_{\mathrm{D}}\iota_V
\end{equation}
inherits the empirical status of the selected $\iota_V$. The calibration of $r_{\mathrm{D}}$ must therefore be interpreted as conditional on
\begin{equation}
    \iota_V = 1\ \mathrm{m}^{-1}
\end{equation}
until the defect-flow boundary-value problem supplies an independently calculated projection. 

The calibration relation may be represented as
\begin{equation}
    \mathcal{F}_{\mathrm{F}} \left( r_{\mathrm{D}}, \iota_V \right) = 0.
\end{equation}

Accordingly,
\begin{equation}
    r_{\mathrm{D}} = r_{\mathrm{D}}(\iota_V).
\end{equation}

A revised microscopic value of $\iota_V$ would generally induce corresponding revisions in
\begin{equation}
    \rho_{\mathrm{IC}}, \qquad K_{\mathrm{D}}, \qquad M_{\mathrm{top}},
\end{equation}
and every subsequent quantity whose calibration depends on the inferred interstitial-condensate impedance.

A complete defect-flow calibration requires an explicit representation of the interstitial FCC geometry rather than a spherical continuum boundary alone. The condensate domain must include the connected tetrahedral and octahedral channels surrounding the defect. The boundary mode must be derived from the displacement or reconfiguration of the localized topological structure, and the response field must be solved with the normal boundary condition
\begin{equation}
    \boldsymbol{\Upsilon}\cdot\hat{\mathbf{n}} = \mathbf{b}_{\partial}\cdot\hat{\mathbf{n}}.
\end{equation}

The projection is then obtained from
\begin{equation}
    A_{\mathrm{imp}}^{(n)} = \int_{\partial\Omega_{\mathrm{def}}} \left| \mathbf{b}_{\partial}\cdot\hat{\mathbf{n}} \right|^2 \mathrm{d}A,
\end{equation}
and
\begin{equation}
    V_{\mathrm{imp}} = \int_{\Omega_{\mathrm{IC}}^{(\mathrm{def})}} |\boldsymbol{\Upsilon}|^2 \mathrm{d}V.
\end{equation}

Convergence must be established with respect to the outer-domain extent, lattice resolution, and numerical discretization. The low-frequency limit must satisfy
\begin{equation}
    \lim_{\omega\rightarrow0} \iota_V^{(\mathrm{def})}(\omega) = \iota_{V,0}^{(\mathrm{def})},
\end{equation}
and the result must be tested across defect geometries and topological branches.

A universal interstitial-condensate impedance density may be defined only when the projection is invariant across a specified stable defect class:
\begin{equation}
    \iota_{V,0}^{(\mathrm{def})} = \iota_{V,0}.
\end{equation}

Under that condition,
\begin{equation}
    \rho_{\mathrm{IC}} = \mu_{\mathrm{D}} \iota_{V,0},
\end{equation}
and
\begin{equation}
    K_{\mathrm{D}} = 4\pi r_{\mathrm{D}}^2 \mu_{\mathrm{D}} \iota_{V,0}.
\end{equation}

If the projection varies among defect classes, the mode-resolved notation must be retained:
\begin{equation}
    \rho_{\mathrm{IC}}^{(\mathrm{imp,def})} = \mu_{\mathrm{D}} \iota_V^{(\mathrm{def})},
\end{equation}
and
\begin{equation}
    K_{\mathrm{D}}^{(\mathrm{def})} = 4\pi r_{\mathrm{D}}^2 \mu_{\mathrm{D}} \iota_V^{(\mathrm{def})} \mathcal{G}_{\mathrm{imp}}.
\end{equation}

The final localized-defect constitutive structure is therefore
\begin{equation}
    \boxed{ \iota_V^{(\mathrm{def})} = \frac{ A_{\mathrm{imp}}^{(n)} }{ V_{\mathrm{imp}} } }
\end{equation}
\begin{equation}
    \boxed{ \rho_{\mathrm{IC}}^{(\mathrm{imp})} = \mu_{\mathrm{D}} \iota_V^{(\mathrm{def})} }
\end{equation}
and
\begin{equation}
    \boxed{ M_{\mathrm{top}} = 4\pi r_{\mathrm{D}}^2 \mu_{\mathrm{D}} \iota_V^{(\mathrm{def})} \mathcal{G}_{\mathrm{imp}} r_{\mathrm{amp}} f_{\varphi}(\Delta\varphi) + \mathcal{O} \left( r_{\mathrm{amp}}^2 \right). }
\end{equation}

The coherent active-vacuum branch and the localized defect branch are thereby separated at the level of their kinetic actions. The coherent branch satisfies
\begin{equation}
    \mathbf{v}_{\mathrm{rel}} = \mathbf{0}
\end{equation}
and supports two finite-frequency translational transverse polarizations subject to the incremental volume constraint. The localized branch satisfies
\begin{equation}
    \mathbf{v}_{\mathrm{rel}} \neq \mathbf{0}
\end{equation}
and acquires a mode-resolved hydrodynamic inertia determined by the defect boundary, phase structure, and induced condensate flow.